\documentclass[11pt,a4paper]{article}
\usepackage[utf8]{inputenc}
\usepackage[russian]{babel}
\usepackage[margin=2cm]{geometry}
\usepackage{listings}
\usepackage{xcolor}
\usepackage{hyperref}
\usepackage{booktabs}
\usepackage{amsmath}
\usepackage{fancyhdr}

\pagestyle{fancy}
\fancyhf{}
\rhead{LIM2 Hardware Simulator Guide}
\lhead{Архитектура, Регистры и Периферия}
\rfoot{Стр. \thepage}

\definecolor{codegreen}{rgb}{0,0.6,0}
\definecolor{codegray}{rgb}{0.5,0.5,0.5}
\definecolor{codepurple}{rgb}{0.58,0,0.82}
\definecolor{backcolour}{rgb}{0.96,0.97,0.98}

\lstdefinestyle{pythonstyle}{
    backgroundcolor=\color{backcolour},
    commentstyle=\color{codegreen},
    keywordstyle=\color{blue}\bfseries,
    numberstyle=\tiny\color{codegray},
    stringstyle=\color{codepurple},
    basicstyle=\ttfamily\small,
    breakatwhitespace=false,
    breaklines=true,
    captionpos=b,
    keepspaces=true,
    numbers=left,
    numbersep=5pt,
    showspaces=false,
    showstringspaces=false,
    showtabs=false,
    tabsize=4
}
\lstset{style=pythonstyle}

\title{\textbf{Полное Архитектурное Руководство по Симулятору LIM2}\\
\large Построчный разбор архитектуры, регистров, флагов, памяти, сигналов и ножек (Pure Python)}
\author{Команда разработки LIM2}
\date{\today}

\begin{document}

\maketitle
\tableofcontents
\newpage

\section{Введение и Философия Проекта}

Симулятор \textbf{LIM2} спроектирован для точной низкоуровневой симуляции процессорной архитектуры и материнской платы с полным сохранением аппаратной гибкости:
\begin{itemize}
    \item \textbf{100\% Чистый Python (Pure Python):} Полное отсутствие внешних тяжелых зависимостей. Все компоненты автономны и прозрачны.
    \item \textbf{Конфигурация платформы в коде:} Вся материнская плата, распределение адресного пространства (Memory Mapping) и список подключенных устройств задаются в коде Python (\texttt{sim.py}). Нет необходимости каждый раз передавать длинные аргументы командной строки (\texttt{-m}, \texttt{-v}).
    \item \textbf{Аппаратные сигнальные линии (Pins):} Полная симуляция линий управления (\texttt{RD}, \texttt{WR}, \texttt{MREQ}, \texttt{IORQ}, \texttt{CS}, \texttt{IRQ}, \texttt{RF} и пользовательских линий).
    \item \textbf{Отставание памяти (Memory Latency \& Wait States):} Возможность задавать задержку выборки ячеек памяти в тактах шины и отрабатывать протокол готовности данных через ножку \texttt{RF} (Ready Flag).
    \item \textbf{Многоколоночный TUI-интерфейс:} Интерактивная визуализация микросхем регистров, потока машинных команд, видеотерминалов и карты распределения памяти в стиле \textit{diskmgmt.msc}.
\end{itemize}

\section{Анатомия Процессора: Построчный Разбор (\texttt{lim2\_cpu.py})}

Процессор наследуется от базового класса \texttt{BaseCPU} и состоит из трех основных методов:
\texttt{setup(self)}, \texttt{on\_start(self)} и \texttt{on\_step(self)}.

\subsection{Создание Регистров (\texttt{setup})}

Каждый регистр создается методом \texttt{self.reg(...)}. Рассмотрим значение каждого аргумента:

\begin{lstlisting}[language=Python, caption={Объявление регистров общего назначения}]
def setup(self):
    for i in range(8):
        self.reg(
            name=f"r{i}",      # 1. Системное имя (доступно как self.r0 .. self.r7)
            bits=16,           # 2. Разрядность в битах (маска 0xFFFF)
            default=0,         # 3. Начальное значение при включении и сбросе
            indexed=True,      # 4. Доступ по номеру: self.get_reg(i) / self.set_reg(i, val)
            title=f"R{i}",     # 5. Подпись на микросхеме в графическом интерфейсе
            role="General",    # 6. Роль регистра
            block="core"       # 7. Функциональный блок процессора
        )
\end{lstlisting}

\begin{itemize}
    \item \texttt{name}: Имя атрибута процессора (\texttt{cpu.r0}, \texttt{cpu.acca}).
    \item \texttt{bits}: Битность регистра. Внутри автоматически рассчитывается маска $\text{mask} = 2^{\text{bits}} - 1$. Запись числа автоматически отсекает лишние старшие биты по маске.
    \item \texttt{indexed=True}: Критически важный параметр! Позволяет обращаться к регистру по номеру: \texttt{self.get\_reg(0)} и \texttt{self.set\_reg(0, val)}. Когда номер регистра декодируется из битов опкода команды, вы напрямую обращаетесь к регистру по его числовому индексу.
\end{itemize}

\subsection{Программный Счетчик (Program Counter, PC)}

Регистр программного счетчика объявляется с флагом \texttt{is\_pc=True}:
\begin{lstlisting}[language=Python]
self.reg("pc", bits=16, default=0x0000, is_pc=True, title="PC")
\end{lstlisting}

\begin{itemize}
    \item \textbf{Ручной инкремент счетчика (\texttt{self.pc += 1}):} Программист сам контролирует продвижение счетчика: \texttt{self.pc += 1} или \texttt{self.pc += 2}. Переполнение разрядной сетки отсекается аппаратной маской автоматически!
    \item \textbf{Чтение байта по адресу PC:} \texttt{val = self.bus.read8(self.pc); self.pc += 1}.
    \item \textbf{Переходы:} Безусловный переход выполняется прямым присваиванием \texttt{self.pc = target\_addr}. Относительный переход --- через \texttt{self.pc += offset}.
\end{itemize}

\subsection{Регистр Флагов (\texttt{FREG})}

Регистр флагов объявляется с параметром \texttt{is\_flag=True} и списком \texttt{flag\_layout}:

\begin{lstlisting}[language=Python, caption={Конфигурация битов регистра флагов FREG}]
flag_layout = [
    # (Имя флага, Полное описание, Номер бита в регистре)
    ("ZF", "Zero", 0),          # Бит 0: Флаг нуля
    ("CF", "Carry", 1),         # Бит 1: Флаг переноса
    ("OF", "Overflow", 2),      # Бит 2: Флаг переполнения
    ("EM", "ErrMath", 3),       # Бит 3: Математическая ошибка
    ("PIE", "ProgInt", 8),      # Бит 8: Разрешение прерываний
    ("MIE", "MasterInt", 9),    # Бит 9: Глобальные прерывания
    ("M24", "24-bit Mode", 11), # Бит 11: Режим 24-битной адресации
]
self.reg("freg", bits=16, default=0, is_flag=True, flag_layout=flag_layout, title="FREG")
\end{lstlisting}

Управление флагами:
\begin{itemize}
    \item \texttt{self.set\_flag("ZF", 1)} --- установить бит флага в 1;
    \item \texttt{self.set\_flag("ZF", 0)} --- сбросить флаг в 0;
    \item \texttt{self.get\_flag("ZF")} --- получить значение флага (0 или 1).
\end{itemize}

\subsection{Выборка и Декодирование Инструкций (\texttt{on\_step})}

\begin{lstlisting}[language=Python, caption={Ручная выборка байтов, инкремент pc += 1 и парсинг в on\_step}]
def on_step(self):
    # 1. Пошаговая выборка двух байтов из шины с ручным инкрементом PC:
    b0 = self.bus.read8(self.pc)
    self.pc += 1

    b1 = self.bus.read8(self.pc)
    self.pc += 1

    # 2. Склейка в 16-битную инструкцию на шине данных databus:
    self.databus = (b0 << 8) | b1

    # 3. Декодирование полей через конструктор парсинга битов:
    opcode, cfg, rd, rs1, rs2 = self.parse("[15-11] [10-8] [7-5] [4-2] [1-0]")
\end{lstlisting}

Значение полей и гибкость парсинга:
\begin{itemize}
    \item \textbf{Диапазоны битов:} \texttt{[15-11]} (5 бит, код операции \texttt{opcode} от 0 до 31).
    \item \textbf{Одиночный бит:} \texttt{[15]} или \texttt{[0]} (извлекает ровно 1 бит, значение 0 или 1). Например: \texttt{is\_imm = self.parse("[15]")}.
    \item \textbf{Парсинг одного поля в переменную:} \texttt{opcode = self.parse("[15-11]")}. Результат ведет себя как полноценное целое число \texttt{int} (поддерживает \texttt{if opcode == 1:}, арифметику, форматирование), а также распаковку через кортеж: \texttt{opcode, = self.parse("[15-11]")}.
    \item \textbf{Смешанные шаблоны:} \texttt{sign, exp, mant = self.parse("[15] [14-10] [9-0]")}.
\end{itemize}

\section{Конфигурация Платформы в Коде (\texttt{sim.py})}

\begin{lstlisting}[language=Python, caption={Объявление устройств в sim.py}]
from bus import MemoryBus
from devices import (
    RAMDevice,
    TextScreenDevice,
    UARTDevice,
    TimerDevice
)

def setup_platform(bus: MemoryBus):
    # 1. Системное ОЗУ: 16 КБ (0x0000 .. 0x3FFF)
    bus.attach(RAMDevice(start=0x0000, end=0x3FFF, name="RAM"))

    # 2. Текстовый видеомонитор 20x4 (0x8000 .. 0x804F)
    bus.attach(TextScreenDevice(start=0x8000, width=20, height=4, name="MONITOR"))

    # 3. Последовательный порт UART (0xF000 .. 0xF003)
    bus.attach(UARTDevice(start=0xF000, name="UART0"))

    # 4. Аппаратный таймер с ножкой IRQ (0xF010 .. 0xF013)
    bus.attach(TimerDevice(start=0xF010, name="TIMER0", irq_pin="IRQ"))
\end{lstlisting}

\section{Аппаратные Ножки, Отставание Памяти (Latency) и Сигнал Готовности (RF)}

В реальных интегральных схемах микросхемы памяти обладают временем доступа. Симулятор \textbf{LIM2} позволяет точно симулировать задержки чтения и протокол рукопожатия (Handshake) с использованием сигнала \texttt{RF} (Ready Flag):

\subsection{Настройка задержки}
\begin{lstlisting}[language=Python]
# ОЗУ с задержкой в 2 такта шины и ножкой готовности RF:
ram = RAMDevice(start=0x0000, end=0x3FFF, latency=2, ready_pin="RF")
bus.attach(ram)
\end{lstlisting}

\subsection{Временная диаграмма цикла шины}
\begin{enumerate}
    \item \textbf{Старт цикла чтения:} Процессор поднимает сигнал \texttt{bus.pins.RD = 1} и выставляет адрес.
    \item \textbf{Захват устройством:} Устройство видит сигнал \texttt{RD}, сбрасывает ножку готовности \texttt{bus.pins.RF = 0} (данные еще не готовы) и запускает таймер задержки.
    \item \textbf{Такты ожидания (Wait States):} Процессор ожидает на шине, пока ножка \texttt{RF == 0}.
    \item \textbf{Готовность данных:} Устройство завершает выборку, выставляет данные на шину и поднимает \texttt{bus.pins.RF = 1}.
    \item \textbf{Захват процессором:} Процессор фиксирует данные и сбрасывает строб \texttt{bus.pins.RD = 0}.
\end{enumerate}

\subsection{Использование в коде процессора}
\begin{lstlisting}[language=Python]
# Чтение с автоматической отработкой тактов ожидания:
data, wait_cycles = self.bus.read_handshake(0x1000)

# Либо пошаговое ожидание ножки RF в цикле инструкций:
waited = self.wait_for_pin("RF")
\end{lstlisting}

\section{Взаимодействие Устройств Друг с Другом и с Процессором}

\begin{itemize}
    \item \textbf{MMIO (Memory-Mapped I/O):} Обращение к регистрам устройств по адресам шины. При чтении/записи по адресу шина автоматически транслирует вызов в соответствующее устройство.
    \item \textbf{Линия прерываний (\texttt{IRQ}):} Таймер при обнулении выставляет \texttt{self.bus.pins.IRQ.set(1)}, а процессор проверяет этот сигнал в начале цикла команды.
    \item \textbf{Слушатели сигналов (\texttt{listen}):} Одно устройство может подписаться на сигнал ножки другого:
    \begin{lstlisting}[language=Python]
    bus.pins.MY_PULSE.listen(lambda val: other_device.on_signal(val))
    \end{lstlisting}
\end{itemize}

\section{Внешние Дополнения: Графический Пиксельный Монитор}

В качестве примера абсолютной расширяемости платформы в папке \texttt{addons/} реализован пиксельный дисплей с графическим окном Tkinter (\texttt{addons/pixel\_monitor.py}):

\begin{lstlisting}[language=Python, caption={Подключение пиксельного монитора в sim.py}]
from addons.pixel_monitor import PixelMonitorDevice

def setup_platform(bus: MemoryBus):
    bus.attach(RAMDevice(0x0000, 0x3FFF, "RAM"))

    # Графический дисплей 64x64 в отдельном окне GUI (масштаб 6x):
    bus.attach(PixelMonitorDevice(
        start=0x4000,
        width=64,
        height=64,
        scale=6,
        mode="rgb332",
        show_gui=True
    ))
\end{lstlisting}

Адрес пикселя $(x, y)$ рассчитывается по формуле:
$$\text{Addr} = \text{start} + (y \times \text{width}) + x$$
Запись цвета процессором в этот адрес мгновенно отображает пиксель в графическом окне ОС.

\section{Инициализация Устройств и Форматы Файлов}

\subsection{Предустановка данных в устройствах (\texttt{default\_data})}
При создании любого устройства в \texttt{sim.py} можно сразу задать его начальное содержимое:
\begin{lstlisting}[language=Python]
# Текстовый экран с приветственной надписью:
bus.attach(TextScreenDevice(0x8000, 20, 4, default_data="LIM2 READY..."))

# ПЗУ с шестнадцатеричным дампом байтов:
bus.attach(ROMDevice(0xF800, 0xFFFF, default_data="08 66 F8 00 12 34"))

# ОЗУ с бинарным файлом ядра:
bus.attach(RAMDevice(0x0000, 0x3FFF, default_data="kernel.bin"))
\end{lstlisting}

\subsection{Поддержка форматов .bin и .hex}
Симулятор автоматически распознает формат входного файла:
\begin{itemize}
    \item \texttt{.bin}: Прямой бинарный образ (байты записываются в память последовательно начиная с адреса \texttt{--org}).
    \item \texttt{.hex}: Шестнадцатеричный текстовый файл или формат Intel HEX (\texttt{:100000...}).
    \item \texttt{.asm / .lasm}: Исходный ассемблерный файл.
\end{itemize}

\section{Режим Непрерывной Сессии (Long-Term Work, \texttt{-L})}

Флаг \texttt{-L} (или \texttt{--long-term [FILE]}) обеспечивает непрерывную долговременную работу:
\begin{lstlisting}[language=bash]
# Запуск с автосохранением в дефолтный дамп lim2_state.dump:
python3 sim.py kernel.bin -i -L

# Запуск с именованным файлом сессии:
python3 sim.py kernel.bin -i -L my_session.dump
\end{lstlisting}

\begin{itemize}
    \item \textbf{При старте:} Если файл снимка существует, восстанавливаются все регистры процессора, счетчик команд \texttt{PC}, такты \texttt{step\_count}, память \texttt{RAM} и состояние всех подключенных устройств.
    \item \textbf{При выходе (\texttt{q} или Ctrl+C):} Симулятор автоматически записывает обновленный снимок машины в указанный файл дампа.
\end{itemize}

\section{Управление Раскладкой и Перенос Блоков по Колонкам}

\subsection{Интерактивное перемещение блоков}
Симулятор разделен на 5 независимых инфо-блоков: \texttt{cpu}, \texttt{code}, \texttt{screen}, \texttt{mem}, \texttt{map}. Любой блок можно динамически переместить в любую колонку:
\begin{lstlisting}[language=bash]
move screen 2      # перенести монитор во вторую колонку
move screen 3      # перенести монитор в третью колонку
move code 1        # перенести дизассемблер в первую колонку
move screen next   # циклический сдвиг в следующую колонку (+)
move screen prev   # сдвиг в предыдущую колонку (-)
cols               # показать текущую раскладку колонок
\end{lstlisting}

\subsection{Автоматический перенос при нехватке высоты экрана (Height Auto-Wrap)}
При небольшой высоте терминала симулятор автоматически отслеживает переполнение колонки. Если суммарная высота блоков превышает размер окна, лишние блоки автоматически перетекают в соседнюю свободную колонку, предотвращая вертикальную прокрутку.
\begin{itemize}
    \item \texttt{wrap on / off}: Включение и выключение автоматического переноса блоков по высоте.
    \item \texttt{height <N|auto>}: Ручная фиксация или автоопределение предельной высоты.
\end{itemize}

\section{Командный Интерфейс Отладчика (\texttt{-i})}

\begin{table}[h!]
\centering
\begin{tabular}{@{}ll@{}}
\toprule
\textbf{Команда} & \textbf{Описание} \\ \midrule
\texttt{<Enter> / s / step} & Выполнить 1 такт/команду \\
\texttt{do <N>} & Выполнить ровно $N$ тактов \\
\texttt{r / run / cont} & Непрерывный запуск симуляции \\
\texttt{delay <ms>} & Установить задержку анимации (мс) \\
\texttt{b <addr>} & Установить/снять точку останова \\
\texttt{r0=0x42 / acca=10} & Изменить значение регистра \\
\texttt{set <addr> <val...>} & Записать байты в память \\
\texttt{write <dev> <off> <val>} & Записать байты или текст в устройство \\
\texttt{save [file]} & Сохранить снимок машины (дамп) в файл \\
\texttt{load [file]} & Загрузить снимок машины или бинарник \\
\texttt{layout [auto|1col|2col|3col]} & Переключить раскладку колонок \\
\texttt{move / cols} & Показать распределение блоков по колонкам \\
\texttt{move <блок> <1|2|3>} & Назначить блок в колонку (напр: \texttt{move screen 2}) \\
\texttt{move <блок> next/prev} & Сдвинуть блок в следующую/предыдущую колонку \\
\texttt{wrap [on|off]} & Автоперенос блоков при нехватке высоты экрана \\
\texttt{height [N|auto]} & Ограничить высоту колонок (в строках) \\
\texttt{view [cpu|code|mem|screen|map]} & Включение/выключение панелей \\
\texttt{map} & Показать карту памяти \textit{diskmgmt} \\
\texttt{reset} & Сброс процессора к начальному состоянию \\
\texttt{q / quit / exit} & Выход (с автосохранением дампа при \texttt{-L}) \\ \bottomrule
\end{tabular}
\caption{Основные команды интерактивного отладчика}
\end{table}

\section{Заключение}

Архитектура \textbf{LIM2} предоставляет разработчику полную свободу: от написания собственного процессора и декодера инструкций до симуляции отставания памяти, аппаратных ножек, межмодульных связей и сложной внешней периферии с наглядной визуализацией в терминале или графических окнах ОС.

\end{document}
